\providecommand{\rank}{\operatorname{r}}
\providecommand{\RR}{\mathbb R}
\providecommand{\CC}{\mathbb C}
\providecommand{\KK}{\mathbb K}

% ---------- Exercise 51 ----------
\setcounter{enumi}{50}
\item {\bfseries\boldmath
For each $k\in\RR$, let
}

\[
A_k=\begin{bmatrix}0&k&2\\k-4&0&0\\0&3&k-1\end{bmatrix},\quad B_k=\begin{bmatrix}1&k&-1\\2&4&-2\\3&-7&k+3\end{bmatrix},\quad C_k=\begin{bmatrix}1&k&k&-2\\0&1&3&1\\1&k&1&-2k\\0&0&k&0\end{bmatrix}.
\]

\begin{enumerate}[label=\alph*)]

\item Compute $\det A_k$, $\det B_k$ and $\det C_k$.

\begin{answer}
\[
 A_k=\begin{bmatrix}0&k&2\\k-4&0&0\\0&3&k-1\end{bmatrix},\quad
 B_k=\begin{bmatrix}1&k&-1\\2&4&-2\\3&-7&k+3\end{bmatrix},\quad
 C_k=\begin{bmatrix}1&k&k&-2\\0&1&3&1\\1&k&1&-2k\\0&0&k&0\end{bmatrix}.
\]

Expansion along row $2$ of $A_k$ gives
\[
 \det A_k=-(k-4)\det\begin{bmatrix}k&2\\3&k-1\end{bmatrix}
 =-(k-4)(k-3)(k+2).
\]
For $B_k$, use $F_2\gets F_2-2F_1$, $F_3\gets F_3-3F_1$:
\[
 \det B_k=\det\begin{bmatrix}1&k&-1\\0&4-2k&0\\0&-7-3k&k+6\end{bmatrix}
 =(4-2k)(k+6).
\]
For $C_k$, use $F_3\gets F_3-F_1$, then expand along row $4$:
\[
 \det C_k=-k\det\begin{bmatrix}1&k&-2\\0&1&1\\0&0&2-2k\end{bmatrix}
 =-k(2-2k)=2k(k-1).
\]
\end{answer}

\item Determine the values of $k$ for which (i) $A_k$ is invertible; (ii) $B_k$ is invertible; (iii) $C_k$ is invertible.

\begin{answer}
The matrices are invertible precisely for
\[
 \begin{array}{ll}
 \text{(i)}\quad k\notin\{-2,3,4\}&\text{for }A_k,\\
 \text{(ii)}\quad k\notin\{-6,2\}&\text{for }B_k,\\
 \text{(iii)}\quad k\notin\{0,1\}&\text{for }C_k.
 \end{array}
\]
\end{answer}

\item Give the values of $k$ for which (i) the linear system $A_kX=0$ has a unique solution; (ii) the linear system $B_kX=0$ has a nonzero solution.

\begin{answer}
(i) $A_kX=0$ has a unique solution exactly when $A_k$ is invertible:
$k\notin\{-2,3,4\}$. (ii) $B_kX=0$ has a nonzero solution exactly when
$B_k$ is singular: $k\in\{-6,2\}$.
\end{answer}

\item Determine the values of $k$ for which $A_k^TB_{2k}$ is invertible.

\begin{answer}
By the determinant product rule,
\[
 \det(A_k^TB_{2k})=\det(A_k)\det(B_{2k}),\qquad
 \det(B_{2k})=(4-4k)(2k+6).
\]
Thus $A_k^TB_{2k}$ is invertible exactly when
$k\notin\{-3,-2,1,3,4\}$.
\end{answer}
\end{enumerate}
